% Môn: Vật lý
% Lớp: 10
% Chương: Chương I. Mở đầu
% Bài: L10C1 Bài 2. Các quy tắc an toàn trong phòng thực hành Vật lí · Xử lí sự cố và sơ cứu trong phòng thực hành
% Số câu: 185
% App đọc file này trực tiếp — sửa rồi Commit trên GitHub, bấm Đồng bộ GitHub .tex

% L10C1 Bài 2. Các quy tắc an toàn trong phòng thực hành Vật lí



% ===== Câu 60 =====
% ID: L10C1B2-137-TN
% Mức: NB
\dangbt{Xử lí sự cố và sơ cứu trong phòng thực hành}
\begin{ex}
% ID: L10C1B2-137-TN
% Mức: NB
Trước khi bắt đầu thí nghiệm liên quan đến xử lí sự cố và sơ cứu, hành động nào an toàn nhất?
\choice
{\True dừng thí nghiệm, cô lập nguồn nguy hiểm, báo giáo viên và sơ cứu đúng nguyên tắc}
{tự xử lí một mình khi chưa bảo đảm an toàn}
{tiếp tục thao tác rồi báo cáo sau}
{nhờ bạn khác làm thay mà không kiểm tra điều kiện an toàn}
\loigiai{
Hành động đúng là “dừng thí nghiệm, cô lập nguồn nguy hiểm, báo giáo viên và sơ cứu đúng nguyên tắc” vì ngăn sự cố lan rộng và hạn chế tổn thương cho nạn nhân. Chọn A.
}
\end{ex}

% ===== Câu 61 =====
% ID: L10C1B2-138-TN
% Mức: NB
\begin{ex}
% ID: L10C1B2-138-TN
% Mức: NB
Khi phát hiện thiết bị có dấu hiệu bất thường liên quan đến xử lí sự cố và sơ cứu, hành động nào an toàn nhất?
\choice
{tự xử lí một mình khi chưa bảo đảm an toàn}
{tiếp tục thao tác rồi báo cáo sau}
{nhờ bạn khác làm thay mà không kiểm tra điều kiện an toàn}
{\True dừng thí nghiệm, cô lập nguồn nguy hiểm, báo giáo viên và sơ cứu đúng nguyên tắc}
\loigiai{
Hành động đúng là “dừng thí nghiệm, cô lập nguồn nguy hiểm, báo giáo viên và sơ cứu đúng nguyên tắc” vì ngăn sự cố lan rộng và hạn chế tổn thương cho nạn nhân. Chọn D.
}
\end{ex}

% ===== Câu 62 =====
% ID: L10C1B2-139-TN
% Mức: NB
\begin{ex}
% ID: L10C1B2-139-TN
% Mức: NB
Trong lúc thay đổi cách bố trí dụng cụ liên quan đến xử lí sự cố và sơ cứu, hành động nào an toàn nhất?
\choice
{tiếp tục thao tác rồi báo cáo sau}
{nhờ bạn khác làm thay mà không kiểm tra điều kiện an toàn}
{\True dừng thí nghiệm, cô lập nguồn nguy hiểm, báo giáo viên và sơ cứu đúng nguyên tắc}
{tự xử lí một mình khi chưa bảo đảm an toàn}
\loigiai{
Hành động đúng là “dừng thí nghiệm, cô lập nguồn nguy hiểm, báo giáo viên và sơ cứu đúng nguyên tắc” vì ngăn sự cố lan rộng và hạn chế tổn thương cho nạn nhân. Chọn C.
}
\end{ex}

% ===== Câu 66 =====
% ID: L10C1B2-140-TN
% Mức: TH
\begin{ex}
% ID: L10C1B2-140-TN
% Mức: TH
Sau khi hoàn thành phép đo liên quan đến xử lí sự cố và sơ cứu, hành động nào an toàn nhất?
\choice
{nhờ bạn khác làm thay mà không kiểm tra điều kiện an toàn}
{\True dừng thí nghiệm, cô lập nguồn nguy hiểm, báo giáo viên và sơ cứu đúng nguyên tắc}
{tự xử lí một mình khi chưa bảo đảm an toàn}
{tiếp tục thao tác rồi báo cáo sau}
\loigiai{
Hành động đúng là “dừng thí nghiệm, cô lập nguồn nguy hiểm, báo giáo viên và sơ cứu đúng nguyên tắc” vì ngăn sự cố lan rộng và hạn chế tổn thương cho nạn nhân. Chọn B.
}
\end{ex}

% ===== Câu 67 =====
% ID: L10C1B2-141-TN
% Mức: TH
\begin{ex}
% ID: L10C1B2-141-TN
% Mức: TH
Khi một bạn chưa đọc hướng dẫn liên quan đến xử lí sự cố và sơ cứu, hành động nào an toàn nhất?
\choice
{\True dừng thí nghiệm, cô lập nguồn nguy hiểm, báo giáo viên và sơ cứu đúng nguyên tắc}
{tự xử lí một mình khi chưa bảo đảm an toàn}
{tiếp tục thao tác rồi báo cáo sau}
{nhờ bạn khác làm thay mà không kiểm tra điều kiện an toàn}
\loigiai{
Hành động đúng là “dừng thí nghiệm, cô lập nguồn nguy hiểm, báo giáo viên và sơ cứu đúng nguyên tắc” vì ngăn sự cố lan rộng và hạn chế tổn thương cho nạn nhân. Chọn A.
}
\end{ex}

% ===== Câu 68 =====
% ID: L10C1B2-142-TN
% Mức: TH
\begin{ex}
% ID: L10C1B2-142-TN
% Mức: TH
Khi khu vực làm việc bị ướt liên quan đến xử lí sự cố và sơ cứu, hành động nào an toàn nhất?
\choice
{tự xử lí một mình khi chưa bảo đảm an toàn}
{tiếp tục thao tác rồi báo cáo sau}
{nhờ bạn khác làm thay mà không kiểm tra điều kiện an toàn}
{\True dừng thí nghiệm, cô lập nguồn nguy hiểm, báo giáo viên và sơ cứu đúng nguyên tắc}
\loigiai{
Hành động đúng là “dừng thí nghiệm, cô lập nguồn nguy hiểm, báo giáo viên và sơ cứu đúng nguyên tắc” vì ngăn sự cố lan rộng và hạn chế tổn thương cho nạn nhân. Chọn D.
}
\end{ex}

% ===== Câu 69 =====
% ID: L10C1B2-143-TN
% Mức: TH
\begin{ex}
% ID: L10C1B2-143-TN
% Mức: TH
Khi cần rời bàn thí nghiệm liên quan đến xử lí sự cố và sơ cứu, hành động nào an toàn nhất?
\choice
{tiếp tục thao tác rồi báo cáo sau}
{nhờ bạn khác làm thay mà không kiểm tra điều kiện an toàn}
{\True dừng thí nghiệm, cô lập nguồn nguy hiểm, báo giáo viên và sơ cứu đúng nguyên tắc}
{tự xử lí một mình khi chưa bảo đảm an toàn}
\loigiai{
Hành động đúng là “dừng thí nghiệm, cô lập nguồn nguy hiểm, báo giáo viên và sơ cứu đúng nguyên tắc” vì ngăn sự cố lan rộng và hạn chế tổn thương cho nạn nhân. Chọn C.
}
\end{ex}

% ===== Câu 70 =====
% ID: L10C1B2-144-TN
% Mức: VD
\begin{ex}
% ID: L10C1B2-144-TN
% Mức: VD
Khi có người không đeo phương tiện bảo hộ liên quan đến xử lí sự cố và sơ cứu, hành động nào an toàn nhất?
\choice
{nhờ bạn khác làm thay mà không kiểm tra điều kiện an toàn}
{\True dừng thí nghiệm, cô lập nguồn nguy hiểm, báo giáo viên và sơ cứu đúng nguyên tắc}
{tự xử lí một mình khi chưa bảo đảm an toàn}
{tiếp tục thao tác rồi báo cáo sau}
\loigiai{
Hành động đúng là “dừng thí nghiệm, cô lập nguồn nguy hiểm, báo giáo viên và sơ cứu đúng nguyên tắc” vì ngăn sự cố lan rộng và hạn chế tổn thương cho nạn nhân. Chọn B.
}
\end{ex}

% ===== Câu 71 =====
% ID: L10C1B2-145-TN
% Mức: VD
\begin{ex}
% ID: L10C1B2-145-TN
% Mức: VD
Khi kết quả đo khác dự kiến liên quan đến xử lí sự cố và sơ cứu, hành động nào an toàn nhất?
\choice
{\True dừng thí nghiệm, cô lập nguồn nguy hiểm, báo giáo viên và sơ cứu đúng nguyên tắc}
{tự xử lí một mình khi chưa bảo đảm an toàn}
{tiếp tục thao tác rồi báo cáo sau}
{nhờ bạn khác làm thay mà không kiểm tra điều kiện an toàn}
\loigiai{
Hành động đúng là “dừng thí nghiệm, cô lập nguồn nguy hiểm, báo giáo viên và sơ cứu đúng nguyên tắc” vì ngăn sự cố lan rộng và hạn chế tổn thương cho nạn nhân. Chọn A.
}
\end{ex}

% ===== Câu 72 =====
% ID: L10C1B2-146-TN
% Mức: VD
\begin{ex}
% ID: L10C1B2-146-TN
% Mức: VD
Khi giáo viên yêu cầu dừng thao tác liên quan đến xử lí sự cố và sơ cứu, hành động nào an toàn nhất?
\choice
{tự xử lí một mình khi chưa bảo đảm an toàn}
{tiếp tục thao tác rồi báo cáo sau}
{nhờ bạn khác làm thay mà không kiểm tra điều kiện an toàn}
{\True dừng thí nghiệm, cô lập nguồn nguy hiểm, báo giáo viên và sơ cứu đúng nguyên tắc}
\loigiai{
Hành động đúng là “dừng thí nghiệm, cô lập nguồn nguy hiểm, báo giáo viên và sơ cứu đúng nguyên tắc” vì ngăn sự cố lan rộng và hạn chế tổn thương cho nạn nhân. Chọn D.
}
\end{ex}

% ===== Câu 133 =====
% ID: L10C1B2-43-TLN
% Mức: NB
\begin{ex}
% ID: L10C1B2-43-TLN
% Mức: NB
Một quy trình về xử lí sự cố và sơ cứu gồm: (1) nhận diện nguy cơ; (2) dừng thí nghiệm, cô lập nguồn nguy hiểm, báo giáo viên và sơ cứu đúng nguyên tắc; (3) tự xử lí một mình khi chưa bảo đảm an toàn; (4) báo giáo viên khi có bất thường; (5) chỉ tiếp tục khi nguy cơ đã được kiểm soát. Có bao nhiêu bước phù hợp quy tắc an toàn? Chỉ ghi số.
\shortans{4}
\loigiai{
Các bước (1), (2), (4), (5) phù hợp; bước (3) không an toàn. Có 4 bước phù hợp.
}
\end{ex}

% ===== Câu 137 =====
% ID: L10C1B2-45-TL
% Mức: VD
\begin{bt}
% ID: L10C1B2-45-TL
% Mức: VD
Xây dựng một bảng kiểm ngắn gồm ít nhất bốn mục cho tình huống khi kết quả đo khác dự kiến liên quan đến xử lí sự cố và sơ cứu.
\loigiai{
Bảng kiểm gồm: đọc hướng dẫn; nhận diện mối nguy; kiểm tra dụng cụ và bảo hộ; thực hiện dừng thí nghiệm, cô lập nguồn nguy hiểm, báo giáo viên và sơ cứu đúng nguyên tắc; theo dõi bất thường; thu dọn và báo cáo sau thí nghiệm.
}
\end{bt}

% ===== Câu 141 =====
% ID: L10C1B2-47-TL
% Mức: VD
\begin{bt}
% ID: L10C1B2-47-TL
% Mức: VD
So sánh biện pháp phòng ngừa và biện pháp ứng phó đối với nguy cơ thuộc nhóm xử lí sự cố và sơ cứu.
\loigiai{
Phòng ngừa diễn ra trước sự cố, tập trung vào kiểm tra, loại bỏ hoặc giảm nguy cơ. Ứng phó diễn ra khi sự cố đã hoặc sắp xảy ra, gồm dừng hoạt động, cô lập nguồn nguy hiểm, báo cáo và sơ cứu. Hai nhóm biện pháp bổ sung cho nhau.
}
\end{bt}

% ===== Câu 154 =====
% ID: L10C1B2-57-DS
% Mức: TH
\begin{ex}
% ID: L10C1B2-57-DS
% Mức: TH
Trong tình huống khi giáo viên yêu cầu dừng thao tác đối với xử lí sự cố và sơ cứu, hãy đánh giá các phát biểu sau.
\choiceTF
{\True Cần dừng thí nghiệm, cô lập nguồn nguy hiểm, báo giáo viên và sơ cứu đúng nguyên tắc.}
{Có thể tự xử lí một mình khi chưa bảo đảm an toàn nếu thao tác diễn ra nhanh.}
{\True Phải báo cho giáo viên hoặc người phụ trách khi phát hiện nguy cơ vượt khả năng kiểm soát.}
{\True Chỉ tiếp tục thí nghiệm sau khi nguyên nhân nguy hiểm đã được loại bỏ hoặc kiểm soát.}
\loigiai{
\begin{itemchoice}
\item (Đúng) Cần dừng thí nghiệm, cô lập nguồn nguy hiểm, báo giáo viên và sơ cứu đúng nguyên tắc. (Vì): vì đây là biện pháp an toàn của dạng xử lí sự cố và sơ cứu; b sai vì hành vi “tự xử lí một mình khi chưa bảo đảm an toàn” vẫn nguy hiểm; c và d đúng do sự cố phải được báo cáo, cô lập và kiểm soát trước khi tiếp tục
\item (Sai) Có thể tự xử lí một mình khi chưa bảo đảm an toàn nếu thao tác diễn ra nhanh.
\item (Đúng) Phải báo cho giáo viên hoặc người phụ trách khi phát hiện nguy cơ vượt khả năng kiểm soát.
\item (Đúng) Chỉ tiếp tục thí nghiệm sau khi nguyên nhân nguy hiểm đã được loại bỏ hoặc kiểm soát.
\end{itemchoice}
}
\end{ex}

% ===== Câu 155 =====
% ID: L10C1B2-58-DS
% Mức: VD
\begin{ex}
% ID: L10C1B2-58-DS
% Mức: VD
Trong tình huống trước khi bắt đầu thí nghiệm đối với xử lí sự cố và sơ cứu, hãy đánh giá các phát biểu sau.
\choiceTF
{\True Cần dừng thí nghiệm, cô lập nguồn nguy hiểm, báo giáo viên và sơ cứu đúng nguyên tắc.}
{Có thể tự xử lí một mình khi chưa bảo đảm an toàn nếu thao tác diễn ra nhanh.}
{\True Phải báo cho giáo viên hoặc người phụ trách khi phát hiện nguy cơ vượt khả năng kiểm soát.}
{\True Chỉ tiếp tục thí nghiệm sau khi nguyên nhân nguy hiểm đã được loại bỏ hoặc kiểm soát.}
\loigiai{
\begin{itemchoice}
\item (Đúng) Cần dừng thí nghiệm, cô lập nguồn nguy hiểm, báo giáo viên và sơ cứu đúng nguyên tắc. (Vì): vì đây là biện pháp an toàn của dạng xử lí sự cố và sơ cứu; b sai vì hành vi “tự xử lí một mình khi chưa bảo đảm an toàn” vẫn nguy hiểm; c và d đúng do sự cố phải được báo cáo, cô lập và kiểm soát trước khi tiếp tục
\item (Sai) Có thể tự xử lí một mình khi chưa bảo đảm an toàn nếu thao tác diễn ra nhanh.
\item (Đúng) Phải báo cho giáo viên hoặc người phụ trách khi phát hiện nguy cơ vượt khả năng kiểm soát.
\item (Đúng) Chỉ tiếp tục thí nghiệm sau khi nguyên nhân nguy hiểm đã được loại bỏ hoặc kiểm soát.
\end{itemchoice}
}
\end{ex}

